\documentclass[10pt,a4paper]{amsart}
\usepackage[margin=19mm]{geometry}
\usepackage{amsmath,amssymb,mathtools}
\usepackage[scr=rsfs,cal=euler]{mathalfa}
\usepackage{xfrac}
\usepackage[unicode,hidelinks,bookmarks=false]{hyperref}
\newcommand{\ie}{i.e.\ }
\newcommand{\cf}{cf.\ }
\newcommand{\resp}{respectively\ }
\newcommand{\Subsection}{\subsection}
\providecommand{\nobreakdash}{\nobreak\hbox{-}}
\setlength{\emergencystretch}{2em}
\multlinegap0pt
\usepackage{bm}
\usepackage{bbm}
\usepackage{dsfont}
\usepackage{xspace}
\usepackage{cancel}
\usepackage{mathtools}
\usepackage{amsthm}
\makeatletter
\@ifundefined{theorem}{\newtheorem{theorem}{Theorem}}{}
\makeatother
\title[Reinhardt domains determined by their endomorphisms]{Reinhardt domains determined by their endomorphisms}
\author{Rafael B. Andrist, W{\l}odzimierz Zwonek}
\date{Source: arXiv:2604.19364v1 (CC BY 4.0)}
\begin{document}
\maketitle

\noindent\emph{Source and licence.} Reinhardt domains determined by their endomorphisms. Version arXiv:2604.19364v1 (CC BY 4.0), distributed under the Creative Commons Attribution 4.0 International licence, as recorded in the arXiv licence metadata for this exact version and shown on the abstract page https://arxiv.org/abs/2604.19364v1. Full rights notice: licenses/arxiv-2604.19364-CC-BY-4.0.txt.\par\smallskip

\noindent\textit{Editorial scope.} Counted unit: the Theorem whose source label is \texttt{thm:counterexample} in the article (source0.tex lines 745--748, source label \texttt{thm:counterexample}), delivered together with the article's own complete original proof of it (source0.tex lines 751--809, one \texttt{proof} environment of 59 lines). No mathematics is written, repaired or re-derived: statement and proof are the article's own text line for line. Required context retained uncounted: none. Every definition, display and label of those lines is retained unchanged. Omitted from this package and not redistributed: 1--744, 749--750, 810--1081. Nothing inside a retained excerpt is omitted or altered: every retained body is delivered exactly as the article's lines are written, so no omission receipt of any kind accompanies this package. Citations: the retained excerpts carry no citation command at all, so no bibliography entry of the article is transcribed and no reference is redistributed. Reference adaptation: none was needed and none was made; every reference inside the delivered unit resolves inside it, so no unresolved reference or citation is produced. Printed slips kept literal: no printed word, symbol, hypothesis or number of the counted statement or of its proof was altered. Layout and class: the article's own class is \texttt{amsart} with packages including amsmath, amssymb, amsthm, amsrefs, inputenc, babel, color; the wrapper is the lane's pinned \texttt{amsart} editorial wrapper at 10pt on A4, which restates the theorem-like environments the retained text needs and adds the article's own macro definitions that the retained text needs (none), with extra wrapper packages (bm, bbm, dsfont, xspace, cancel, mathtools). The delivered numbering is the unit's own. Assets: no retained span contains a figure environment, a graphics-inclusion command, a PDF-page inclusion, a table or a TikZ picture; no image file is copied into this package, the package declares no asset dependency and no image file is redistributed with it. Licence: version arXiv:2604.19364v1 of this article is distributed under the Creative Commons Attribution 4.0 International licence, as recorded in the arXiv licence metadata for this exact version and shown on the abstract page https://arxiv.org/abs/2604.19364v1. A full rights notice is delivered at licenses/arxiv-2604.19364-CC-BY-4.0.txt. Characters: every retained excerpt is pure ASCII, so the delivered engine, XeLaTeX with its default Latin Modern text font, reports no missing character.
\begin{theorem}
\label{thm:counterexample}
There exist uncountably many holomorphically inequivalent non-pseudoconvex Reinhardt domains $D\subset\mathbb C^2$ such that $\operatorname{End}(D)$ consists only of rotations, i.e.\ mappings $D\owns z\to(\omega_1z_1,\omega_2z_2)\in D$ for $|\omega_j|=1$, and of constant functions.
\end{theorem}

\begin{proof} We start by recalling the endomorphisms of the domain $D_{\gamma}$, where $\gamma>1$ is not algebraic. It consists of the functions $(0,h)$, $(h,0)$ and functions of the form 
\begin{equation}\beta\cdot z^n\cdot p_{\gamma}\circ h,
\end{equation}
where $h\in\mathcal O(D_{\gamma})$, $n=0,1,2,\ldots$, $0<|\beta_1||\beta_2|^{\gamma}\leq 1$, where the last inequality is strict when $n=0$.

Now we define the domain $D$ as follows.
\begin{equation}
    D:=D_{\gamma}\setminus\left((\mathbb C\setminus\mathbb D)\times\{0\}\cup \{0\}\times(\mathbb C\setminus\mathbb D)\cup I_1\cup I_2\right),
    \end{equation}
    where $I_1:=\partial\triangle(0,2)\times \overline{\triangle}(0,1/4^{1/\gamma})$, $I_2:=(\mathbb C\setminus\triangle(0,1/2^{\gamma+1}))\times\partial\triangle(0,2)$.

The set $D$ is the union of $D_{\gamma}^*$ and of two unit discs lying in the axes and with two slits deleted. 

In our case, the essential property is that the logarithmic images of all surfaces $V_r(\gamma)\cap D$ are equal to $V_r(\gamma)$ with one torus removed (for all $r\in(0,1)$, $r\neq 1/2$) or two tori removed (for $r=1/2$).

Note that any holomorphic function on $D$ extends to a holomorphic function on $D_{\gamma}$, so we shall work further with the functions defined on the latter domains. Moreover, one may verify that any endomorphism of $D$ extends to an endomorphism of $D_{\gamma}$. Note also that $D$ is dense in $D_{\gamma}$. Our purpose is to show that if the endomorphism of $D_{\gamma}$ maps $D$ to $D$ then it must be of the form as claimed in the theorem.

We consider first endomorphism of the form $(h,0)$ (and analoguously, $(0,h)$). Note that on $D$ the values of $h$ are bounded by one, so $h$ is bounded by one on $D_{\gamma}$. But any bounded holomorphic function on $D_{\gamma}$ is constant.

Consider the other situation. Note that these endomorphisms (being also endomorphisms of $D_{\gamma}^*$) have the following property: They map the real analytic hypersurface $V_r(\gamma)$ with $0<r<1$ to the surface of the same type $V_s(\gamma)$. 

We prove the following fact that is crucial in the proof of the theorem.

\bigskip

{\bf Claim.} The function $h$ appearing in the formula of an endomorphism must be constant. 

\begin{proof}[Proof of Claim] To show the Claim, it is sufficient to show that for some (equivalently, any) $\beta\in D_{\gamma}^*$ the function
\begin{equation}
 \Phi \colon \mathbb C\owns\lambda \mapsto h(\beta\cdot p_{\gamma}(\lambda))\in\mathbb C
\end{equation}
is constant. To see that this is sufficient, assume $\Phi$ is constant. Then $h$ is constant on $\beta\cdot p_{\gamma}(\mathbb C)$ which, by the irrationality of $\gamma$, is dense in
$V_{|\beta_1||\beta_2|^{\gamma}}(\gamma)$. Consequently, the function $h$ is constant on $V_{|\beta_1||\beta_2|^{\gamma}}(\gamma)$. Then the identity principle shows that $h$ is constant.

Suppose that there is an endomorphism $\Phi(z)=\beta\cdot z^n\cdot p_{\gamma}(h(z))$, $z\in D$ with $h$ being not constant. Choose $r_0\in(0,1)$, $r_0\neq 1/2$. Denote the torus from $V_{r_0}(\gamma)$ not belonging to $D$ by $T$. We claim that $\Phi$ maps $T$ to a torus. To prove that, it is sufficient to show that the mapping $\Phi_0$, where $\Phi_0(z):=\beta\cdot p_{\gamma}(h(z))$, would map $T$ to a torus. If this were the case then the mapping 
\begin{equation}
    D_{\gamma}\owns z \mapsto \log |\beta_2 \exp(h(z)|=\log|\beta_2|+\operatorname{Re}(h(z))
\end{equation}
would be a pluriharmonic mapping that would be constant on the torus $T$. The identity principle would imply that $h$ is constant.

The form of $\Phi$ and the geometry of $D$ let us conclude that the preimages of tori (one or two) from a surface $V_s(\gamma)$ not belonging to the domain $D$ when intersected with some (any) surface $V_r(\gamma)$ must be contained in tori (one or two) from a surface.

Note that by Picard's theorem, for any $a\in T$ the set $\Phi_0(a\cdot p_{\gamma}(\mathbb C))$ touches both tori $T_1$ and $T_2$ that are deleted from the domain $D$. Moreover, to establish the fact that $\Phi_0(T)$ is a torus, it is sufficient to show that if $\Phi_0(a)\in T_1$ then for all $\lambda\in\mathbb C$ with $\operatorname{Re}(\lambda)=0$ we get that $\Phi_0(a\cdot p_{\gamma}(\lambda))\in T_1$. Suppose the opposite. Then, as the mapping
\begin{equation}
    \lambda\mapsto h(a\cdot p_{\gamma}(\lambda))
\end{equation}
is not constant, by the openness of holomorphic functions, we would get the existence of $\lambda$ arbitrarily close to zero with $\operatorname{Re}(\lambda)\neq 0$ and such that $\operatorname{Re}(h(a\cdot p_{\gamma}(\lambda)))=0$, which implies that for $z_0:=a\cdot p_{\gamma}(\lambda)\in V_{r_0}(\gamma)\setminus T$ we get that $\Phi_0(z_0)\in T_1$  -- a contradiction. And this finishes the proof of Claim. 
\end{proof}

%\bigskip

By Claim we know that $\Phi_0(z)=(\beta_1z_1^n,\beta_2z_2^n)$. 

 From the geometry of $D$ (the slits!) we easily get that $n=0$ or $n=1$ (and then $|\beta_1|=|\beta_2|=1$). To see this, one may look at the geometry of $\log D$
 %\begin{equation}
  %   \log D:=\{x\in\mathbb R^n:(e^{x_1},\ldots,e^{x_n})\in D\}
 %\end
 that has to be preserved by the affine mapping $\log D\owns x \mapsto (\log|\beta_1|+n x_1,\log|\beta_2|+nx_2)\in\log D$.
\end{proof}

\end{document}
